\documentclass[11pt,a4paper]{article}
\usepackage{../../common/preamble}

\begin{document}

\begin{center}
    {\LARGE \textbf{Принципы устройства компьютеров}}\\[0.2cm]
    \rule{\textwidth}{1pt}
\end{center}

\section{Принципы фон Неймана}

В 1946 году была запущена первая крупная электронная вычислительная машина ENIAC. Несмотря на применение радиоламп, процесс смены программы на ENIAC был механическим: инженеры переставляли штекерные перемычки на коммутационных панелях и выставляли тысячи тумблеров. Машина не была универсальным автоматом в современном понимании.

Осмыслив опыт создания ENIAC, математик Джон фон Нейман совместно с физиками Артуром Берксом и Германом Голдстайном в июне 1946 года подготовили технический отчет \textit{\guillemotleft Предварительное рассмотрение логической конструкции электронного вычислительного устройства\guillemotright}. Сформулированные в нем принципы определили архитектуру ЭВМ на десятилетия вперед.

Обычно выделяют следюущие наиболее важные идеи этой работы:
\begin{enumerate}
    \item Состав основных компонентов вычислительной машины
    \item Принцип двоичного кодирования
    \item Принцип адресности памяти
    \item Принцип иерархической организации памяти
    \item Принцип хранимой программы
    \item Принцип программного управления
\end{enumerate}

\section{Общие принципы}

\subsection{Состав основных компонентов ЭВМ}
В классической схеме фон Неймана все периферийные потоки данных и обмен с памятью физически проходили через \textbf{центральный процессор}:

\begin{center}
\begin{tikzpicture}[scale=0.85, every node/.style={transform shape}]
    \draw[thick, fill=black!5] (4, 1.5) rectangle (7.5, 3.5);
    \node[align=center] at (5.75, 2.5) {\textbf{Процессор}\\\footnotesize (АЛУ, УУ)};

    \draw[thick, fill=white] (0, 1.5) rectangle (2.5, 3.5);
    \node[align=center] at (1.25, 2.5) {\textbf{Внутренняя}\\\textbf{память}};

    \draw[thick, fill=white] (9, 1.5) rectangle (11.5, 3.5);
    \node[align=center] at (10.25, 2.5) {\textbf{Внешняя}\\\textbf{память}};

    \draw[thick, fill=white] (4.5, 4.5) rectangle (7, 5.7);
    \node[align=center] at (5.75, 5.1) {\textbf{Устройства}\\\textbf{ввода}};

    \draw[thick, fill=white] (4.5, -0.7) rectangle (7, 0.5);
    \node[align=center] at (5.75, -0.1) {\textbf{Устройства}\\\textbf{вывода}};

    \draw[thick, <->, Stealth-Stealth] (2.5, 2.5) -- (4, 2.5);
    \draw[thick, <->, Stealth-Stealth] (7.5, 2.5) -- (9, 2.5);
    \draw[thick, ->, -Stealth] (5.75, 4.5) -- (5.75, 3.5);
    \draw[thick, ->, -Stealth] (5.75, 1.5) -- (5.75, 0.5);
\end{tikzpicture}
\end{center}

\begin{enumerate}
    \item \textbf{Арифметико-логическое устройство (АЛУ):} аппаратный вычислитель, выполняющий элементарные арифметические (сложение, вычитание) и логические (конъюнкция, дизъюнкция, инверсия, битовые сдвиги) операции.
    \item \textbf{Устройство управления (УУ):} блок синхронизации и координации, управляющий выполнением машинных команд. В современных ЭВМ АЛУ и УУ объединены на одном полупроводниковом кристалле и образуют \textbf{микропроцессор}.
    \item \textbf{Память:} разделена на \textit{внутреннюю} (для хранения данных и команд непосредственно во время вычислений) и \textit{внешнюю} (для долговременного хранения между сеансами).
    \item \textbf{Устройства ввода/вывода:} согласующие преобразователи физических воздействий в электрические цифровые коды и обратно.
\end{enumerate}

\subsection{Принцип двоичного кодирования}
Вся обрабатываемая и хранимая информация кодируется в двоичной системе. Полупроводниковые транзисторы в ключевом режиме (состояния насыщения и отсечки) обладают максимальной помехоустойчивостью и технологической простотой.

\begin{infobox}{Троичная ЭВМ \guillemotleft Сетунь\guillemotright}
В 1959 году в МГУ под руководством Н.~П.~Брусенцова была спроектирована малая ЭВМ \guillemotleft Сетунь\guillemotright, работавшая в симметричной троичной системе (разряды $-1, 0, +1$). 
Однако проект остался оригинальным эпизодом: промышленность уже была стандартизирована под двоичные элементы.
\end{infobox}

\section{Принципы организации памяти}

\subsection{Принцип адресности памяти}
Оперативная память представляет собой одномерное упорядоченное пространство дискретных ячеек.

\begin{definitionbox}{Ячейка и адрес памяти}
\begin{itemize}
    \item \textbf{Ячейка памяти} --- это минимально возможный считываемый из памяти объем данных. Невозможно аппаратно прочитать меньшее количество бит (а тем более отдельный изолированный бит). В современных ЭВМ стандартная ячейка равна 1 байту (8 бит).
    \item \textbf{Адрес ячейки памяти} --- ее уникальный порядковый номер. Нумерация адресов начинается с нуля.
\end{itemize}
\end{definitionbox}

\subsection{Эволюция ячейки}
В первых ЭВМ память была числовой: одна ячейка имела большую разрядность и вмещала ровно одно вещественное число или одну инструкцию (что упрощало устройство памяти с инженерной точки зрения). Позднее стало нужно обрабатывать символьный текст. Хранить один символ алфавита (требующий 8 бит) в такой ячейке было расточительно. Инженеры ввели стандарт \textbf{байтовой адресации}:

\begin{center}
\begin{tikzpicture}[scale=0.85]
    \node[align=center, font=\bfseries] at (1.5, 4.5) {а) Ранние ЭВМ};
    \foreach \y/\addr in {0/205, 0.6/204, 1.2/203, 1.8/202, 2.4/201, 3.0/200} {
        \draw[thick] (0, \y) rectangle (4, \y+0.6);
        \node[left=4pt] at (0, \y+0.3) {\texttt{\addr}};
    }
    \node at (2.0, 1.5) {\textit{Числа / Команды}};

    \node[align=center, font=\bfseries] at (8.5, 4.5) {б) Современные ЭВМ};
    \foreach \y/\addr in {0.6/207, 1.2/206, 1.8/205, 2.4/204} {
        \draw[thick] (6, \y) rectangle (7.25, \y+0.6);
        \node[left=4pt] at (6, \y+0.3) {\texttt{\addr}};
    }

    \foreach \y/\addr in {0/208, 3.0/200} {
        \draw[thick] (6, \y) rectangle (11, \y+0.6);
        \node[left=4pt] at (6, \y+0.3) {\texttt{\addr}};
        \draw[dashed] (7.25, \y) -- (7.25, \y+0.6);
        \draw[dashed] (8.5, \y) -- (8.5, \y+0.6);
        \draw[dashed] (9.75, \y) -- (9.75, \y+0.6);
    }

    \foreach \x/\addr in {7.25/201, 8.5/202, 9.75/203} {
        \node[above=2pt] at (\x, 3.6) {\tiny \texttt{\addr}};
    }

    \node at (8.5, 3.3) {\footnotesize 32-битное число (4 байта)};
    \node at (8.75, 1.8) {\footnotesize Символы текста};
    \node at (8.5, 0.3) {\footnotesize 32-битное число};
\end{tikzpicture}
\end{center}

\textbf{Примечание}: если число занимает несколько байт (например, 32-битный \texttt{int} требует 4 последовательные ячейки \texttt{200, 201, 202, 203}), то адресом всего числа считается наименьший из адресов ячеек (в данном примере --- адрес \texttt{200}).

\subsection{Произвольный доступ к памяти}
Память, построенная по принципу адресности, обладает свойством \textbf{произвольного доступа} (\textit{Random Access}): время чтения или записи информации в любую ячейку не зависит от ее адреса (в отличие от, например, магнитных лент с последовательным доступом, где ленту требовалось физически перематывать).

Исторически сложилось, что в русском языке термину RAM ставится в полное соответствие термин ОЗУ (оперативное запоминающее устройство), хотя фактически к RAM относится не только ОЗУ:

\begin{itemize}
    \item \textbf{ОЗУ (Оперативное запоминающее устройство):} Энергозависимая память с возможностью многократного быстрого чтения и записи.
    \item \textbf{ПЗУ (Постоянное запоминающее устройство, ROM, Read Only Memory):} Энергонезависимая память только для чтения. Хранит микрокод первичного запуска, выполняемый аппаратно сразу после включения питания, когда в ОЗУ еще нет данных.
\end{itemize}

\subsection{Принцип иерархической организации памяти}
К памяти предъявляются 2 противоречивых требования: пользователю требуется \textbf{максимальный объем} памяти при \textbf{минимальной задержке доступа} и низкой стоимости. Очевидно, физически ни одно устройство не удовлетворяет обоим требованиям сразу: любое существенное увеличение количества памяти приводит к замедлению работы, а любое существенное ускорение - к увеличению стоимости и, следовательно, уменьшению объема. Таким образом, чтобы бороться с этим противоречием, память организуют иерархически:
$$\text{Регистры ЦП} \xrightarrow{\text{быстрее, но меньше}} \text{Кэш-память} \to \text{ОЗУ} \to \text{Внешние накопители (SSD/HDD)}$$

\section{Принципы, связанные с исполнением программ}

\subsection{Принцип хранимой программы}

Как было сказано ранее, из-за своей механической природы настройка программы для ENIAC могла занимать несколько дней. Чтобы исправить это неудобство, был предложен принцип хранимой программы: все команды представляются в виде двоичного кода и записываются заранее на некоторые носители, с которых по необходимости данная программа полностью считывается в память компьютера.

\subsection{Принцип однородности памяти}

Поскольку команды и сами данные по принципу хранимой программы стали по форме представления одинаковыми, появилась возможность хранить их в единой памяти. Это свойство назвали принципом однородности памяти.

\begin{questionbox}{Главное следствие принципа однородности}
Какое колоссальное преимущество дает программистам тот факт, что процессор хранит команды и данные в общем пространстве в одинаковом формате?
\tcblower
\textbf{Ответ:} Команды программы могут формироваться как результат вычислений другой программы. Благодаря этому принципу стало возможным создание \textbf{компиляторов} и \textbf{трансляторов}: программа считывает исходный текст на языке C++/Python как обычные текстовые данные и записывает в память машинные команды, готовые к непосредственному выполнению процессором.
\end{questionbox}

\subsection{Гарвардская архитектура}

В Гарвардском университете под руководством профессора Говарда Эйкена по заказу флота США была построена электромеханическая машина Harvard Mark I:
\begin{itemize}
    \item Команды считывались с 24-канальной перфоленты (память команд).
    \item Числа хранились на электромеханических счетных колесах и реле (память данных).
\end{itemize}

Шины считывания команд и шины передачи данных были физически изолированы. Так возникла альтернативная архитектурная ветвь, в которой не выполняется принцип однородности памяти.

\begin{center}
\begin{tikzpicture}[scale=1.0, every node/.style={transform shape}]
    \node[font=\bfseries] at (3, 3.4) {Архитектура фон Неймана};
    \draw[thick, fill=black!5] (0, 0.8) rectangle (2.5, 2.6);
    \node[align=center] at (1.25, 1.7) {\textbf{ЦПУ}};
    \draw[thick, fill=black!8] (4.5, 0.8) rectangle (7, 2.6);
    \node[align=center] at (5.75, 1.7) {\textbf{Общее ОЗУ}\\\footnotesize (Код + Данные)};
    \draw[thick, <->, Stealth-Stealth] (2.5, 1.7) -- (4.5, 1.7) node[midway, above] {\footnotesize Общая шина};

    \draw[dotted, thick] (7.8, 0) -- (7.8, 3.6);

    \node[font=\bfseries] at (12, 3.4) {Гарвардская архитектура};
    \draw[thick, fill=black!5] (8.5, 0.8) rectangle (11, 2.6);
    \node[align=center] at (9.75, 1.7) {\textbf{ЦПУ}};
    \draw[thick, fill=black!8] (13.5, 1.9) rectangle (16.5, 3.0);
    \node[align=center] at (15.0, 2.45) {\footnotesize Память программ};
    \draw[thick, fill=black!8] (13.5, 0.4) rectangle (16.5, 1.5);
    \node[align=center] at (15.0, 0.95) {\footnotesize Память данных};
    \draw[thick, <->, Stealth-Stealth] (11, 2.45) -- (13.5, 2.45) node[midway, above] {\tiny Шина команд};
    \draw[thick, <->, Stealth-Stealth] (11, 0.95) -- (13.5, 0.95) node[midway, above] {\tiny Шина данных};
\end{tikzpicture}
\end{center}

Сегодня гарвардская архитектура в основном используется для создания микроконтроллеров.

\subsection{Принцип программного управления}

\subsubsection{Счетчик адреса команд и машинный цикл}

Автоматическое исполнение программ реализуется при помощи специального регистра процессора --- \textbf{счетчика адреса команд} (далее этот регистр будет обозначаться \textbf{PC}, но в разных архитектурах он может называться по-разному). 
В нем всегда хранится адрес команды, которая должна быть выполнена следующей.

\begin{definitionbox}{Основной алгоритм работы процессора}
\begin{enumerate}
    \item \textbf{Выборка команды:} из ячейки оперативной памяти по адресу из регистра $PC$ считывается очередная команда и сохраняется в регистре команд ($IR$).
    \item \textbf{Инкремент счетчика:} значение $PC$ увеличивается на размер считанной команды так, чтобы указывать на следующую по порядку инструкцию.
    \item \textbf{Исполнение:} УУ декодирует код операции, считывает операнды, а АЛУ производит вычисление. Результат записывается в регистр или память.
    \item Переход к шагу 1.
\end{enumerate}
\end{definitionbox}

\textbf{Команды перехода (ветвления и циклы):}  
Естественный последовательный ход выполнения нарушается командами безусловного (\texttt{JMP}) или условного (\texttt{JZ}, \texttt{JNZ}) перехода. Если условие перехода истинно, процессор на шаге 3 записывает в счетчик $PC$ адрес перехода, и следующая выборка происходит из новой ветки кода.

\textbf{Как запускается компьютер при подаче питания?}
В первых ЭВМ оператор вручную тумблерами заносил в регистр $PC$ начальный адрес. В современных компьютерах при подаче питания на линии аппаратного сброса (\textit{Reset}) в регистр $PC$ аппаратно прошивается константный стартовый адрес ПЗУ (\textit{Reset Vector}). Там располагается программа тестирования оборудования и первичной инициализации (BIOS/UEFI), которая находит накопитель, загружает загрузчик операционной системы в ОЗУ и передает управление ему.


\subsubsection{Конвейерный метод обработки данных}

В классическом цикле фазы выборки команды из ОЗУ и ее исполнения в АЛУ происходят по очереди. При этом АЛУ простаивает во время выборки из памяти, а блок выборки простаивает во время работы АЛУ.

\begin{center}
\begin{tikzpicture}[scale=0.8]
    \node at (-1.5, 2.5) {\footnotesize Команда 1:};
    \node at (-1.5, 1.5) {\footnotesize Команда 2:};
    \node at (-1.5, 0.5) {\footnotesize Команда 3:};
    
    \foreach \t in {1, 2, 3, 4, 5} {
        \node at (\t*2 - 1, 3.3) {\footnotesize Такт \t};
    }
    
    \draw[thick, fill=black!10] (0, 2) rectangle (2, 3) node[midway] {\footnotesize Выборка};
    \draw[thick, fill=black!20] (2, 2) rectangle (4, 3) node[midway] {\footnotesize Декод};
    \draw[thick, fill=black!30] (4, 2) rectangle (6, 3) node[midway] {\footnotesize Расчет};
    
    \draw[thick, fill=black!10] (2, 1) rectangle (4, 2) node[midway] {\footnotesize Выборка};
    \draw[thick, fill=black!20] (4, 1) rectangle (6, 2) node[midway] {\footnotesize Декод};
    \draw[thick, fill=black!30] (6, 1) rectangle (8, 2) node[midway] {\footnotesize Расчет};
    
    \draw[thick, fill=black!10] (4, 0) rectangle (6, 1) node[midway] {\footnotesize Выборка};
    \draw[thick, fill=black!20] (6, 0) rectangle (8, 1) node[midway] {\footnotesize Декод};
    \draw[thick, fill=black!30] (8, 0) rectangle (10, 1) node[midway] {\footnotesize Расчет};
\end{tikzpicture}
\end{center}

\textit{Суть конвейера:} разбиение выполнения команды на независимые подоперации. На каждом такте конвейер выдает результат одной готовой инструкции, повышая быстродействие в несколько раз.

\textit{Проблема конвейера:} при условных переходах предвыбранные команды из следующей строки кода оказываются ненужными --- конвейер приходится сбрасывать, теряя такты.

\end{document}